% ============================================================
% Capítulo 7 — Conclusiones
% ============================================================

\chapter{Conclusiones}
\label{cap:conclusiones}

% ------------------------------------------------------------
\section{Conclusiones}
% ------------------------------------------------------------

\paragraph{Fundamentos (O1).} La homología persistente da un lenguaje
preciso para los errores que más pesan en la segmentación de
estructuras finas. Un filamento partido, dos filamentos fusionados o una
componente espuria son cambios en el número de Betti $\beta_0$, y su
robustez frente al umbral se lee en el diagrama de persistencia.
Topograph traduce esa idea en una pérdida que actúa solo sobre las
regiones críticas, con un coste casi lineal.

\paragraph{Protocolo (O2).} Todas las decisiones se toman en una mitad
de la validación y la cifra se reporta en la otra, con la métrica
oficial y una partición por días. Esta separación ha resultado
necesaria: la variación entre semillas, de hasta $0{,}024$ en PQ, es del
mismo orden que las diferencias entre pérdidas, y el \emph{bootstrap}
por días no la recoge. Sin esta precaución, varias de las comparaciones
habrían parecido concluyentes.

\paragraph{Diagnóstico (O3).} Leído sobre el diagrama $H_0$, el
post-proceso clásico es un filtro sobre las componentes del conjunto de
supernivel que solo mira su tamaño. Su efecto es grande: el área mínima
sube el PQ de la baseline de $0{,}311$ a $0{,}373$, porque la mayoría de
los falsos positivos son componentes pequeñas. El umbral, por sí solo,
no aporta nada.

\paragraph{Pérdidas topológicas (O4).} Sin post-proceso, Topograph
mejora a la baseline con las dos semillas probadas, y es la única pérdida
con una mejora significativa. clDice reduce la fragmentación, coherente
con que premia la continuidad. Pero la mejora de Topograph viene sobre
todo de eliminar componentes espurias, que es lo mismo que hace el
filtro de área. Con el post-proceso ajustado, ninguna pérdida supera a
la baseline. La conclusión es que, en este problema, el sesgo
topológico introducido en el entrenamiento y el filtro clásico aplicado
después corrigen el mismo tipo de error, y sus mejoras no se suman.
\textbf{[completar: barbas y ajuste de hiperparámetros.]}

% ------------------------------------------------------------
\section{Limitaciones}
% ------------------------------------------------------------

\begin{itemize}
  \item \emph{Una ejecución por pérdida.} Solo la baseline y Topograph
  se han entrenado con dos semillas. Dada la variación observada entre
  semillas, las diferencias pequeñas entre pérdidas no pueden
  confirmarse.
  \item \emph{Resolución.} El modelo trabaja a $512\times512$, un cuarto
  de la resolución original en cada eje. Para las barbas esto apenas
  importa (la anotación reducida conserva el 99,7\,\% de su esqueleto),
  pero limita la precisión con que se delimita el contorno.
  \item \emph{Configuración fija.} Las comparaciones de pérdidas usan los
  mismos hiperparámetros, sin aumentos de datos y con el peso de
  Topograph fijado a priori. El ajuste de la
  Sección~\ref{sec:res-hp} cubre solo tres modelos.
  \item \emph{Una sola arquitectura}, la U-Net con ResNet-34.
\end{itemize}

% ------------------------------------------------------------
\section{Trabajo futuro}
% ------------------------------------------------------------

\begin{itemize}
  \item \textbf{Post-proceso por persistencia (E2).} Es la continuación
  más directa. Si Topograph y el filtro de área corrigen el mismo error,
  un filtro que decida sobre la persistencia de cada componente, y no
  sobre su área, podría conservar los filamentos pequeños pero robustos
  que el filtro de área elimina. La fusión por proximidad del collado
  atacaría además los fragmentados recuperables, que ningún método
  probado corrige.
  \item \textbf{Estrategias de etiquetas (E4).} Comparar el
  entrenamiento con cada anotación por separado con la media de las
  anotaciones y con los filamentos de la mayoría.
  \item \textbf{Más semillas por pérdida}, para separar las diferencias
  reales del ruido del entrenamiento.
  \item \textbf{Pérdidas centradas en las barbas.} Como la resolución
  no es el límite, mejorar las barbas pasa por cambiar lo que aprende
  la red, por ejemplo dando más peso a la zona de barbas en la pérdida.
  \item \textbf{Otras pérdidas topológicas}, como Betti
  matching~\cite{stucki2023}, que garantiza la correspondencia espacial
  de las estructuras a mayor coste computacional.
\end{itemize}
